\documentclass{article}

\title{Shakespearean Sonnets}
\author{Anony Mouse}
\date{}

\begin{document}

\maketitle

William Shakespeare was the author of many plays, including tragedies like \textit{Romeo and Juliet}, comedies like \textit{A Midsummer Night's Dream}, and histories like \textit{Richard II}. Or was he? 

There is ample controversy as to whether Shakespeare himself wrote these works, or if he served as the stand-in for another author such as Christopher Marlowe or Francis Bacon. We will likely never know who the true author of Shakespeare's works was, but ``William Shakespeare,'' whoever he \textit{actually} was, also wrote a series of sonnets. One of the most famous of these is Sonnet 18:

\begin{verse}

Shall I compare thee to a summer’s day? \\
Thou art more lovely and more temperate: \\
Rough winds do shake the darling buds of May, \\
And summer’s lease hath all too short a date; \\
Sometime too hot the eye of heaven shines, \\
And often is his gold complexion dimm'd; \\
And every fair from fair sometime declines, \\
By chance or nature’s changing course untrimm'd; \\
But thy eternal summer shall not fade, \\
Nor lose possession of that fair thou ow’st; \\
Nor shall death brag thou wander’st in his shade, \\
When in eternal lines to time thou grow’st:  \\
So long as men can breathe or eyes can see, \\
So long lives this, and this gives life to thee. \\

\end{verse} 

As you can see from this example, Shakespearean sonnets consist of 14 lines of iambic pentameter. They typically have three quatrains followed by a couplet, follwing the rhyme scheme \textbf{ABAB CDCD EFEF GG}. 

The verse environment indents like the quote environment, but it is ragged right instead of fully justified.

\end{document} 





